\documentclass[UTF8,a4paper,11pt,fontset=fandol]{ctexart}
\usepackage[margin=25mm]{geometry}
\usepackage{amsmath,amssymb,amsthm,mathtools,bm}
\usepackage{booktabs,longtable,array,enumitem}
\usepackage[numbers,sort&compress]{natbib}
\usepackage{xurl}
\usepackage[colorlinks=true,linkcolor=blue,citecolor=blue,urlcolor=blue]{hyperref}
\usepackage{microtype}
\providecommand{\doi}[1]{\href{https://doi.org/#1}{\nolinkurl{#1}}}
\setlength{\emergencystretch}{3em}
\setlength{\parskip}{0.25em}
\setlist{nosep}
\numberwithin{equation}{section}
\newtheorem{theorem}{定理}[section]
\newtheorem{proposition}[theorem]{命题}
\newtheorem{lemma}[theorem]{引理}
\newtheorem{corollary}[theorem]{推论}
\theoremstyle{definition}
\newtheorem{definition}[theorem]{定义}
\newtheorem{assumption}[theorem]{假设}
\theoremstyle{remark}
\newtheorem{remark}[theorem]{说明}
\newtheorem{counterexample}[theorem]{反例}
\DeclareMathOperator{\Fix}{Fix}
\DeclareMathOperator{\Stab}{Stab}
\DeclareMathOperator{\rank}{rank}
\DeclareMathOperator{\tr}{tr}
\DeclareMathOperator{\im}{im}
\DeclareMathOperator{\diag}{diag}
\DeclareMathOperator{\Span}{span}
\newcommand{\R}{\mathbb R}
\newcommand{\C}{\mathbb C}
\newcommand{\Z}{\mathbb Z}
\newcommand{\Q}{\mathbb Q}
\newcommand{\T}{\mathsf T}
\newcommand{\I}{\mathrm I}
\newcommand{\eps}{\varepsilon}
\newcommand{\norm}[1]{\left\lVert #1\right\rVert}
\newcommand{\code}[1]{\texttt{\detokenize{#1}}}
\newcommand{\scope}{本文是项目技术论证，未经独立同行评审。定理证明针对明列的数学假设；
它们不自动证明某个二进制、数据库、浮点容差或官网版本满足这些假设。}

\title{CRIS 中缺少直接实现依据的计算环节：\\数学与晶体学补充论证}
\author{CRIS 项目理论资料整理}
\date{2026年10月15日}
\begin{document}
\maketitle
\begin{abstract}
本文针对文献给出一般原理、但不能直接认证 CRIS 具体实现约定的环节，提供可检验的补充证明。
结果包括有限群精确稳定子的充要条件与有理见证、有限超胞的平移特征标、
斜晶格最近像的有限搜索界、模式子空间与列身份的区别、正定度量匹配的唯一解、
应变右作用、离散谐波的独立性，以及截断搜索能够声称的完备性。
全文区分文献结论、在本文给出的推导和仍待实证的实现条件。
未找到直接信源并不允许臆造来源，也不意味着可以把未经证明的软件行为改名为定理。
\end{abstract}
\noindent\textbf{关键词：}各向同性子群；精确稳定子；周期晶格；模式归一化；有限应变；可验证计算。
\par\scope
\clearpage
{\small\tableofcontents}
\clearpage

\section{问题、证据等级与约定}
ISODISTORT 的一般理论来自空间群表示、序参量和模式投影
\citep{Campbell2006,Stokes1991,Hatch2002,Stokes2017}。
但是，这些论文没有逐行认证本仓库的匹配器、标签重排、容差、缓存或导出代码。
“没有找到直接实现依据”在本文指：检索到的原始文献不足以确定某一具体程序分支的合同，
并非断言整个数学问题从未被研究。
本文只证明明确陈述的合同，不以特定材料的已知答案代替一般论证。

证明使用三种标记：\textbf{已证明}表示在列出的假设下有完整推导；
\textbf{条件性支持}表示算法成立但输入表示、完整性或坐标约定仍须核验；
\textbf{未认证}表示文献与本次推导均不能认证实际软件行为。
这些标记不代表对实测精度的评级。
文献检索与源码对应关系由同目录 \code{coverage.json} 和 \code{audit/} 保存。

除专门说明外，晶格基矩阵 $A$ 的列为笛卡尔直接基，位置为 $r=Ax$，$x$ 是分数列坐标。
$\{R|t\}$ 作用为 $x\mapsto Rx+t$；乘法 $gh$ 先作用 $h$。
倒易分数坐标 $k$ 的平移相位采用 $\exp(2\pi i k^\T n)$。
另一种 Fourier 符号会改变相位正负，但必须贯穿生产者和消费者。
有限群表示 $D:G\to GL(V)$ 默认是真正的同态，不把近似闭包或未经验证的矩阵集合当作群。

\section{精确稳定子判据及可构造的有理见证}\label{sec:stabilizer}
\begin{theorem}[Reynolds 投影]
设 $G$ 有限，$H\leq G$，$D$ 是实或复表示。定义
\[
 P_H=\frac1{|H|}\sum_{h\in H}D(h).
\]
则 $P_H^2=P_H$ 且 $\im P_H=\Fix(H)$，其中
$\Fix(H)=\{v:D(h)v=v,\ h\in H\}$。
若内积满足 $D(g)^*WD(g)=W>0$，则 $P_H$ 是该内积下的正交投影。
\end{theorem}
\begin{proof}
双重求和中每个 $k\in H$ 恰出现 $|H|$ 次，故
$P_H^2=|H|^{-2}\sum_{h,j}D(hj)=P_H$。
对任意 $h_0\in H$，左乘置换求和项，$D(h_0)P_H=P_H$，
故像包含于不变子空间。反之 $v\in\Fix(H)$ 时 $P_Hv=v$。
最后 $W^{-1}P_H^*W=|H|^{-1}\sum_hD(h)^{-1}=P_H$。
\end{proof}
这是群平均原理的自足证明；其晶体学应用背景见
\citet{Hatch2002,Stokes2017}。不能对无限空间群直接写有限平均；
公度情形必须先给出使表示下降的正规平移子群。

\begin{theorem}[非零精确稳定子的充要条件]\label{thm:exact}
以下将 $V$ 视为实向量空间；复表示可先按实部、虚部实化。
令 $V_H=\Fix(H)$，$K_H=\{g\in G:D(g)v=v\ \forall v\in V_H\}$。
存在非零 $v$ 满足 $\Stab_G(v)=H$，当且仅当
\[
 \dim V_H>0,\qquad K_H=H.
\]
\end{theorem}
\begin{proof}
若这样的 $v$ 存在，则 $H\subseteq K_H\subseteq\Stab(v)=H$。
反之，对 $g\notin H$，$W_g=V_H\cap\ker(D(g)-\I)$ 是 $V_H$ 的真线性子空间。
有限多个真线性子空间不能覆盖一个正维实向量空间：为每个 $W_g$ 选一个在
$V_H$ 上不恒为零、却在 $W_g$ 上为零的线性函数 $\ell_g$。
非零多项式 $\prod_g\ell_g$ 不可能在整个 $\R^d$ 上恒为零，
这一事实可由变量个数归纳及一元多项式只有有限根证明。
故可选 $v$ 同时避开所有 $W_g$，并避开原点；其稳定子恰是 $H$。
若 $H=G$，直接取 $V_H$ 中任一非零向量。
\end{proof}

\begin{proposition}[有限整数参数给出有理见证]\label{prop:witness}
若全部表示矩阵有理，$v_1,\ldots,v_d$ 是 $V_H$ 的有理基，且定理
\ref{thm:exact} 的条件满足，令
$q(t)=\sum_{j=1}^d t^{j-1}v_j$。
当 $H\ne G$ 时，在整数
$0\leq t\leq (|G|-|H|)(d-1)$ 中至少有一个给出精确稳定子 $H$。
\end{proposition}
\begin{proof}
对每个 $g\notin H$，$(D(g)-\I)q(t)$ 至少有一个分量是不恒零、
次数不超过 $d-1$ 的多项式，否则 $g$ 固定全部基向量。
各自至多排除 $d-1$ 个实数；有限集合的并至多排除
$(|G|-|H|)(d-1)$ 个整数。给定区间多一个整数，必有可用者。
由于 $v_1$ 的系数为 $1$，$q(t)$ 永不为零。
\end{proof}
这是确定性构造，不需要随机采样碰到“足够一般”的 OPD。
$\Fix(H)\ne0$ 单独只证明“至少具有 $H$”，不能证明“恰好具有 $H$”。
当 $\Fix(H)=0$ 时，只剩零畸变；其稳定子是整个 $G$，不是任意候选 $H$。
对应代码为 \code{backend/method3/inverse_landau.py}。
有理算术可保证这些有限线性运算精确，但不能保证输入的 irrep 矩阵来自正确数据库。

\begin{proposition}[耦合序参量]
在直和表示 $V=\bigoplus_i V_i$ 中，
$\Stab((v_i)_i)=\bigcap_i\Stab(v_i)$。
\end{proposition}
\begin{proof}
一个算符固定直和向量，当且仅当逐分量固定；两边条件完全相同。
\end{proof}
若把若干分量通过非单射映射加到同一物理位移空间，可能发生抵消，此时不能直接套用。
表示允许某种畸变也不表示其自由能一定下降。
主序参量的选择不必唯一；这在逆 Landau 论文中已被指出\citep{Hatch2002,Coupled1991}。

\section{平移格、波矢与有限折叠}\label{sec:lattice}
\begin{theorem}[超胞与波矢的精确相容性]
在母相原胞坐标中，子平移格为 $S\Z^3$，
$S\in M_3(\Z)$、$\det S\ne0$。
波矢 $k\in\R^3/\Z^3$ 的模式对全部子平移不变，当且仅当
\[
 S^\T k\in\Z^3.
\]
这样的不同波矢共有 $|\det S|$ 个。
\end{theorem}
\begin{proof}
对子平移 $Sn$ 的相位为 $\exp(2\pi i k^\T Sn)$。
它对所有 $n\in\Z^3$ 等于 $1$，等价于逐基向量检验
$S^\T k\in\Z^3$。
映射 $\ell\mapsto S^{-\T}\ell$ 给出
$\Z^3/S^\T\Z^3$ 与这些波矢的双射。
对任意整数非奇异矩阵 $J$，伴随矩阵恒等式给出
$(\det J)\Z^3\subseteq J\Z^3$，所以商群有限，记其阶为 $N$。
取 $N$ 个整数余类代表 $z_i$，则 $\bigcup_i(z_i+[0,1)^3)$
是子格 $J\Z^3$ 的一个基本域，体积为 $N$；$J[0,1)^3$ 是另一个基本域，
体积为 $|\det J|$。两个有界基本域通过与格平移的交集分割、平移重组，
保持体积，边界零测度，故 $N=|\det J|$。取 $J=S^\T$ 即得计数。
\end{proof}
实现也可使用 Smith 标准形 $USV=\diag(s_1,s_2,s_3)$；使用某个库的分解结果时必须独立检查
$USV$、整数性、$\det U=\det V=\pm1$，不能只信返回的对角数字。
这与逆 Landau 的平移相位约束及公度算法一致
\citep{Hatch2002,Stokes2017}。

\begin{proposition}[有理相位阶数]
若原胞基上的 $k_j=a_j/b_j$ 已约分，则特征标 $\chi_k$ 的阶为
$\operatorname{lcm}(b_1,b_2,b_3)$。
对多行调制矩阵 $Q$，允许平移为
$L_Q=\{n\in\Z^3:Qn\in\Z^d\}$；其指数等于
$\Z^3\to(\Q/\Z)^d,\ n\mapsto Qn$ 的像大小。
\end{proposition}
\begin{proof}
$m\chi_k=0$ 等价于 $mk_j$ 全为整数，即每个 $b_j\mid m$。
第二个结论由第一同构定理 $\Z^3/L_Q\simeq\im Q$ 得出；
有理 $Q$ 的分母公倍数使其像有限。
\end{proof}
多个波矢的联合指数一般不能由逐行阶数的简单乘积代替，因为像可能相关。
原胞条件也不能由惯用胞的三个平移替代：
例如体心平移 $(1/2,1/2,1/2)$ 对惯用胞 $k=(1,0,0)$ 的相位是 $-1$，
所以这个向量不是母相倒易格的零余类。
对应代码为 \code{backend/utils/lattice.py} 和
\code{backend/method2/superspace.py}；所有“整数性”须在实际原胞基中核验。

\begin{theorem}[谐波别名和实模式维数]\label{thm:harmonics}
设有限平移商群 $F$ 上的特征标 $\chi$ 阶为 $m$。
$\chi^a=\chi^b$ 当且仅当 $a\equiv b\pmod m$。
不同特征标正交；$\chi^a$ 与 $\chi^{-a}$ 形成一个二维实
cosine/sine 子空间，除非 $2a\equiv0\pmod m$，此时该特征标为实值、只有一维。
\end{theorem}
\begin{proof}
第一式由阶的定义。对非平凡特征标 $\psi$，
选 $t$ 使 $\psi(t)\ne1$；平移置换给出
$\sum_x\psi(x)=\psi(t)\sum_x\psi(x)$，所以和为零。
应用于 $\chi^a\overline{\chi^b}$ 得正交性。
当共轭特征标不相同时，其实部和虚部线性无关；
相同时，特征标值为 $\pm1$，虚部为零。
\end{proof}
本结论证明离散周期上的别名和维数，不证明网站参数
\code{nmod} 的具体 UI 含义，也不允许把无理波矢自动替换为有限超胞。
真实非公度结构需要超空间模型，公度近似的误差必须另外给出。

\section{周期匹配的有限、度量正确的搜索}\label{sec:matching}
\begin{theorem}[最近晶格像的完备有限界]
设 $A$ 非奇异、$d\in\R^3$，目标为
$\min_{n\in\Z^3}\norm{A(d-n)}_2$。
任选 $n_0$ 得上界 $b=\norm{A(d-n_0)}$，令 $\sigma_{\min}>0$ 为 $A$ 的最小奇异值。
每个更优解都满足
\[
 |d_j-n_j|\leq b/\sigma_{\min}\quad(j=1,2,3).
\]
故枚举该有限整数盒并计算笛卡尔距离可获得全局最近像。
\end{theorem}
\begin{proof}
$\norm{Ay}\geq\sigma_{\min}\norm y$；
若某坐标超过该界，则 $\norm{A(d-n)}>b$，不可能优于已知候选。
闭有界盒含有限整数点，最小值存在。
\end{proof}
这是完备算法的界，不是声称仓库任何现有局部搜索已经实现了该界。
\begin{counterexample}[分数坐标逐分量舍入并不充分]
取二维斜晶格列基
$A=\left(\begin{smallmatrix}1&0.9\\0&0.1\end{smallmatrix}\right)$，
$d=(0.49,0.49)^\T$。
逐分量舍入给 $n=(0,0)$，距离约 $0.9323$；
而 $n=(1,0)$ 的距离约 $0.0846$。可沿第三方向加独立基，得到三维反例。
\end{counterexample}

\begin{proposition}[固定合同下的全局配对]
固定母子相对应、物种集合和晶格度量后，用精确周期距离构造代价
$c_{ij}$，并在物种相容的双射上求
$\min_\pi\sum_i c_{i,\pi(i)}$，可获得该代价意义下的全局最优配对。
\end{proposition}
\begin{proof}
允许配对构成有限集合；精确遍历即证明最小值存在且可求。
使用指派算法时，其输出须满足双射并达到该离散优化的最优值；
算法正确性不能由各原子分别取最近邻推出。
\end{proof}
若有并列最优解、原点自由度或同种原子交换，最优距离不提供唯一物理身份。
群作用要求置换还满足复合律，否则构造出来的“表示”可能不是表示。
相关数值对称性背景见 \citet{Findsym2005,Spglib2024}；
本次没有审计其所有容差分支。

\section{模式基、归一化与来源身份}\label{sec:modes}
\begin{theorem}[带度量的分解与残差]
设 $W>0$，$B\in\R^{n\times r}$ 满列秩。对位移 $u$，
\[
 a=(B^\T WB)^{-1}B^\T Wu,\qquad
 \Pi_B=B(B^\T WB)^{-1}B^\T W
\]
给出唯一最小二乘系数，且 $B^\T W(u-Ba)=0$。
\end{theorem}
\begin{proof}
对 $\norm{u-Ba}_W^2$ 求导得到正规方程。
$B^\T WB$ 正定，所以解唯一。代回得残差正交及 $\Pi_B^2=\Pi_B$。
若 $B$ 不满列秩，唯一的是正交投影而非所有系数；须使用明确的广义逆合同。
\end{proof}

\begin{proposition}[子空间相同不等于逐列相同]\label{prop:columns}
满列秩 $B,C$ 满足 $\im B=\im C$ 当且仅当存在可逆矩阵 $T$ 使 $C=BT$。
此时 $a_C=T^{-1}a_B$。只有进一步确定 $T$ 的标签、尺度、符号和排序，才能转移逐模式幅度。
\end{proposition}
\begin{proof}
若空间相同，$C$ 各列在 $B$ 中有唯一坐标，组成 $T$；
等秩保证 $T$ 可逆。反向立即成立。由 $Ba_B=BTa_C$ 及单射性得系数关系。
\end{proof}
例如 $B=\I_2$，$C$ 为 $45^\circ$ 旋转矩阵；两投影均是单位矩阵，但同一向量的系数不同。
因此“同 irrep、同维数、同子空间”不足以证明官网的 \code{a1,a2,...} 标签一致。
需要原子排列、坐标基、相位原点、Wyckoff 轨道和列变换的来源证据。
这是 \code{backend/method2/superspace.py} 与导出合同必须额外检查的事项。

\begin{proposition}[原胞/超胞幅度换算]
若归一化按原始未应变几何的物理度量进行，$A_s^2$ 是含 $N$ 个母相原胞的
子相原胞中的总平方贡献，$A_p^2$ 定义为每母相原胞的平均平方贡献，
则 $A_s^2=N A_p^2$，从而
\[
 A_p=A_s/\sqrt N=A_s\sqrt{V_p/V_s}.
\]
\end{proposition}
\begin{proof}
将子相原胞的总平方贡献除以所包含的母相原胞数即得平均贡献；
取平方根并保留所选模式的符号。非零 k 模式在各母相原胞中不必相同，
所以不能要求每个局部贡献都等于这个平均值。
$N$ 为原始参考原胞体积比，不能把中心化惯用胞的体积比直接当成原子重数。
\end{proof}
这一合同见官方帮助的完整模式说明\citep{OfficialHelp}。
斜晶格分数位移 $f$ 的范数为 $f^\T A^\T Af$，不是 $f^\T f$。
若 $\hat b_i=n_i b_i$，实际位移为 $\sum_i A_{s,i}n_i b_i$；
物理幅度、原始线性系数、\code{normfactor} 和最大原子位移是不同量。
不同物理量的范数具有不同单位，不能将位移与磁矩的平方直接相加。
应变为宏观变量，不能自动应用上述重复原胞规则。

\section{斜晶格中度量匹配的唯一解}\label{sec:strain}
令代码中的行晶格矩阵为 $P=A^\T$，行超胞变换为 $B$，子晶格为 $D$。
定义
\[
 G=PP^\T,\qquad H=B^{-1}DD^\T B^{-\T}.
\]
假设 $P,B,D$ 非奇异，故 $G,H>0$。
\begin{theorem}[对称正定乘子的唯一性]\label{thm:spd}
方程 $MGM=H$ 有唯一对称正定解
\[
 M=G^{-1/2}\bigl(G^{1/2}HG^{1/2}\bigr)^{1/2}G^{-1/2}.
\]
\end{theorem}
\begin{proof}
置 $X=G^{1/2}MG^{1/2}>0$，则 $X^2=G^{1/2}HG^{1/2}>0$。
实对称正定矩阵具有唯一正定平方根：
谱分解构造平方根；任何正定平方根与其平方交换，逐特征子空间上的
正特征值必须是原特征值的正平方根，故唯一。
代回获得所述公式。直接相乘验证原方程。
\end{proof}
进一步，$Y=BMP$ 满足 $YY^\T=DD^\T$，
因此 $O=Y^{-1}D$ 满足 $OO^\T=\I$。
若 $\det D$ 与 $\det Y$ 同号，则 $O\in SO(3)$；否则只有正交等价而不是纯旋转。
数值矩阵平方根和极分解的背景见 \citet{Higham1986}。

\begin{remark}[已证明的内容与尚未证明的官网对应]
本定理严格支持 \code{backend/method4/strain.py} 中的正定度量方程。
它不自行证明官网在所有斜晶格情形都使用这个对称规范。
$M$ 是所选分数基坐标中的对称乘子；
$AMA^{-1}$ 一般不对称，故 $M-\I$ 不能不经变换就称为笛卡尔有限应变张量。
换非正交基后“欧氏对称”也不自动保持。其与官网具体版本的对应必须核对公式或输出。
官方版本记录明确指出 2013 年更改过应变的归一化与施加方式\citep{OfficialVersions}。
\end{remark}

\begin{theorem}[线性化的可逆性与右作用]\label{thm:right}
对正定 $G$，$\mathcal L_G(E)=EG+GE$ 在对称矩阵空间上可逆。
令 $\mathcal C_R(S)=R^\T SR$，
则 $\mathcal A_R=\mathcal L_G^{-1}\mathcal C_R\mathcal L_G$ 满足
\[
 \mathcal A_R\mathcal A_S=\mathcal A_{SR}.
\]
\end{theorem}
\begin{proof}
在 $G=U\diag(\lambda_i)U^\T$ 的基中，
$\mathcal L_G$ 将第 $ij$ 个对称分量乘以 $\lambda_i+\lambda_j>0$，故可逆。
又 $\mathcal C_R(\mathcal C_S(X))=(SR)^\T X(SR)$，
相似传输保持此复合次序。
\end{proof}
该形式是右作用；若接口要求左作用，应使用 $\mathcal A_{R^{-1}}$。
群不变空间和对整个群求平均不受取逆重排影响，但不能未经说明把矩阵乘积方向反过来。
工程剪切坐标为
$q=(E_{11},E_{22},E_{33},2E_{23},2E_{13},2E_{12})$；
Frobenius 范数在该坐标上的权重是
$\diag(1,1,1,\tfrac12,\tfrac12,\tfrac12)$，并非单位矩阵。
该线性化证明不能用于认证有限应变下任意线性叠加具有完全相同的非线性群作用。

\section{非公度微观模式：从函数作用到有限线性方程}\label{sec:micro}
2017 年各向同性子群论文的宏观序参量构造不能单独确定所有微观列
\citep{Stokes2017}。这里补出通用构造；它与超空间及磁调制约束的思路一致
\citep{Stokes2013,Magnetic2012}，但不替代官网的列命名规范。

\begin{assumption}[明确可检查的输入]
设 $X$ 是含原子站点、晶格平移及内部环面坐标的构形标签空间。
给定群在其上的左作用 $z\mapsto gz$；原子置换、跨胞平移及内部相移
均已包括在这个作用中。场的张量作用为一个真正的表示 $T(g)$
（占位标量、极矢量或含时间反演的轴矢量）。
所保留的有限 Fourier/站点函数空间 $V$ 在全部所需操作下封闭。
\end{assumption}
\begin{theorem}[微观表示及其固定模式]
定义 $(D(g)f)(z)=T(g)f(g^{-1}z)$。则 $D(gh)=D(g)D(h)$。
若 $H$ 由有限集合 $h_1,\ldots,h_s$ 生成，则
\[
 \Fix_V(H)=\ker\begin{bmatrix}D(h_1)-\I\\ \vdots\\D(h_s)-\I\end{bmatrix}.
\]
此式不要求 $H$ 或其相位像有限。
\end{theorem}
\begin{proof}
$D(g)D(h)f(z)=T(g)T(h)f(h^{-1}g^{-1}z)=D(gh)f(z)$。
若各生成元固定 $f$，其逆也固定，任意有限乘积均固定；反向由包含关系立即成立。
由于 $V$ 不变，选定基后全部操作由有限矩阵精确表示，方程的核就是所求空间。
\end{proof}
例如内部仿射作用 $\xi\mapsto S\xi+\delta$ 的主动拉回满足
\[
 e^{2\pi i m^\T S^{-1}(\xi-\delta)}
 =e^{-2\pi i(S^{-\T}m)^\T\delta}
  e^{2\pi i(S^{-\T}m)^\T\xi}.
\]
所以它将索引 $m$ 送至 $S^{-\T}m$ 并乘一个确定相位。
站点相关相位必须由 $g^{-1}z$ 中实际的跨胞/内部位移计算，不能假定所有站点同相。
所选谐波集合必须在这些索引置换和复共轭下封闭；任意截断若不封闭，
截断矩阵的乘法就未必还是表示。数值 Fourier 采样也不能代替这一代数检查。

\begin{proposition}[指定 irrep 的微观嵌入]
给定有限维不可约表示 $\Delta$ 和上述 $D$，求矩阵 $C$ 满足
\[
 D(g_i)C=C\Delta(g_i)
\]
对生成元全部成立，即可求出全部等变嵌入的线性空间。
非零解为单射；若 $V$ 已完全可约，则所有这些解的像张成
$V$ 中与 $\Delta$ 同型的全部分量。
\end{proposition}
\begin{proof}
生成元等式通过乘法和取逆延伸到整个群。
$\ker C$ 是 $\Delta$ 的不变子空间，所以非零解的核为零。
在完全可约直和中，每个同型不可约分量的包含映射都给出一个解；
而不同型分量上的投影与 $C$ 的复合由 Schur 引理为零。
故解的所有像恰好张成目标等型分量。实表示的交换子代数可能不是 $\R$；
不能把解空间维数直接当作不可约重数。
\end{proof}
有限站点置换、正交张量作用和环面仿射变换保持相应的求和/$L^2$ 内积；
在不变有限子空间上，其表示是酉的。因而不变子空间的正交补仍不变，
递归分解证明这里的完全可约性，无须对无限群错误地作有限平均。
选择原点、站点基、内坐标相位和张量坐标后，本节给出严格的微观模式构造。
若其中任何输入尚未知，则只能得到条件性证明；官网的排序、符号及重复实例标签
仍需单独的来源合同。

\section{误差、穷尽性与不能由论文推出的结论}\label{sec:validation}
\begin{proposition}[固定模型下的扰动界]
若 $\widetilde B=W^{1/2}B$ 满列秩，且只改变观测 $u$，则
$\norm{\delta a}_2\leq
\norm{W^{1/2}\delta u}_2/\sigma_{\min}(\widetilde B)$。
\end{proposition}
\begin{proof}
$\delta a=\widetilde B^+W^{1/2}\delta u$；
奇异值分解给 $\norm{\widetilde B^+}_2=1/\sigma_{\min}(\widetilde B)$。
\end{proof}
这是固定配对、固定模式基的界。原点、匹配、空间群或秩发生离散变化时，不能继续使用同一线性误差模型。
很小的残差既不保证正确身份，也不保证系数稳定。
容差应由输入不确定度与尺度说明，不能用一个裸常数证明所有晶体正确。

\begin{proposition}[有界搜索的条件性完备性]
若目标集合 $S$ 有限，枚举器恰好覆盖 $S$，判定谓词 $P$ 精确，
则“枚举后保留 $P$ 为真者”恰好得到 $\{x\in S:P(x)\}$。
若枚举器只访问 $S'\subsetneq S$，无论已访问案例多么准确，均不能由此推出全局完备。
\end{proposition}
\begin{proof}
前一句由集合推导式直接成立；后一句可取未访问元素 $s\in S\setminus S'$ 且 $P(s)$ 为真作为反例。
\end{proof}
应用时必须写明 $S$：所选 k 集、irrep 组合、平移指数、原点类与物理类型。
对连续 k 空间、无上界超胞或资源限时的搜索，“跑完一批”不是穷尽性证明。
同理，CIF 成功写入、GUI 进程成功启动或与一例官网结果相同均不蕴含晶体学正确。

\section{证明义务清单与仍待核验事项}\label{sec:obligations}
本文已经证明的是数学合同，而非以下尚未完成的事实：
\begin{enumerate}
\item 所有 \code{iso/findsym/smodes} 二进制和数据库条目的实现正确性与完整性。
未公开实现不能仅凭接口输出得到逐行证明。
\item 全部斜晶格、极端条件数及有限应变下，代码规范与当前官网规范逐项相同。
需要版本固定的独立输出、坐标合同及数值稳定性分析。
\item 任意重复 irrep 实例的列排序、符号、原点相位和微观来源与官网一一对应。
命题 \ref{prop:columns} 说明子空间等价为何不足。
\item 真实非公度结构在有限谐波截断后的误差界与完备性。
定理 \ref{thm:harmonics} 不能替代对无限 Fourier 展开的收敛假设。
\item 全产品域的输入解析、退化原子匹配、所有导出器、GUI 和并发行为。
这些是实现验收问题，必须以代码审计和可追溯测试补齐。
\end{enumerate}
因此，对任何尚无证明或证据的能力，正确交付状态是“未认证”，而不是制造一个无条件结论。
进一步工作应以本资料库的证明义务为验收条件，避免把本次文献整理当成整套科学软件的正式认证。

\section{共享晶格与站点映射合同}\label{sec:shared-mapping}
本节为本项目推导，使用晶格与坐标变换的一般理论\citep{ITCGeometry,ITA}，
不声称是官网实现的逐行证明。为对应代码，本节单独采用分数行向量：
直接基矩阵 $A_r=A^\T$ 的行为笛卡尔基，子胞 $A'_r=BA_r$。

\begin{proposition}[心化指数与唯一站点对应]
设 $P$ 的行是以惯用基表示的原胞基，则惯用胞中原胞个数为
$m=1/|\det P|$。对未畸变子胞站点行坐标 $x'$, 若母胞代表点为 $x_i$，
其对应关系为 $x'B=x_i+n$, $n\in\mathbb Z^3$。
在给定数值误差界 $\varepsilon\geq0$ 下，只有物种及占据相同且满足
\[
 \|(x'B-x_i-n)A_r\|_2\leq\varepsilon,\qquad n\in\mathbb Z^3
\]
的母胞代表与整数平移二元组 $(i,n)$ 恰为一个时，才可据此唯一选定站点与平移。
\end{proposition}
\begin{proof}
原胞体积为 $|\det(PA_r)|=|\det P|\,|\det A_r|$，故体积比即所述指数。
由 $x'A'_r=(x'B)A_r$ 得坐标关系；周期等价位置相差整数平移，
每个候选像的笛卡尔误差即不等式左端。无可行二元组时输入不满足映射合同，多个二元组时
该误差界不能证明唯一性；两者均不能通过取首行或默认索引零来补齐。
即使只有一个母胞站点，它也可能有两个满足误差界的整数像，因此必须同时检查平移。
\end{proof}
这里 $\varepsilon$ 的单位为 \AA；用分数坐标欧氏距离排序在一般斜晶格中不等价。
实现用晶格库枚举误差球内的全部周期像，再按物种与占据筛选；恰有一个二元组才通过。
该判据相对于数值误差界是保守的：即使最近者严格更近，只要两个像都在误差球内，
仍拒绝宣称唯一性。验证应覆盖站点重排、斜晶格、超胞整数平移以及单站点多像反例。

\section{可复核的计算路径}
后续开发应逐项保存：实际基矩阵与原胞转换、整数格指数、群乘法与表示同态检查、
固定空间的秩、精确稳定子、原子双射、模式 Gram 矩阵和列变换、
未应变归一化、应变度量残差、实现及外部引擎的签名。
本文件内的代数证明不依赖任何特定材料名称、候选编号或预先保存的答案。
\par\textbf{结论。}在明确的有限群、精确晶格和正定度量假设下，上述核心合同可严格推出。
其在 CRIS 和官网中的适用性必须通过同一假设清单逐项确认，不能从证明文本越级推断。

\bibliographystyle{unsrtnat}
\bibliography{references}
\end{document}
